% GENERATED FILE. Edit canonical lessons, then run npm run generate:books.
% canonical-source-hash: fnv1a64:df7d825cb4d187f0
% canonical-lessons: FR-C454-falaise, FR-C454-plateau, FR-C454-sommet, FR-C454-col, FR-C454-versant

\chapter{Le relief}
\label{ch:fr-le-relief}
\hlchaptermodality{454}

\begin{chapteropening}
\textbf{By the end of this chapter:} \emph{I can describe high ground: the cliff, the plateau, the summit, the mountain pass and the side of a mountain.}

Complete the last lesson of chapter 454: I can describe high ground: the cliff, the plateau, the summit, the mountain pass and the side of a mountain.
\end{chapteropening}

\section[la falaise]{la falaise — the cliff}

\label{lesson:FR-C454-falaise}



\begin{quote}
Before the new one: say the French for the cardboard box, then the French for the registered letter.
\end{quote}

\subsection*{You'll want to know: la falaise}
\textbf{la falaise} — \textquotedblleft{}the cliff\textquotedblright{}.

\begin{quote}
Du haut de la falaise, on voit toute la baie.
\end{quote}

\begin{quote}
\emph{From the top of the cliff the whole bay is in view.}
\end{quote}

Usually over the sea.

\begin{cousinweb}
\emph{falaise} comes from Old French \emph{faleise}, from Frankish \emph{falisa} 'rock, cliff' (EtymDB).
\end{cousinweb}

\begin{grammarlens}[title={What you've built}]
One more word for telling, explaining and describing at B1.
\end{grammarlens}

\subsection*{Your turn}
\begin{itemize}
\raggedright
  \item \emph{Say it:} \emph{la falaise}
  \item \emph{Say it:} \emph{la falaise}, once more
  \item \emph{Say it:} a sentence of your own with \emph{la falaise}
  \item \emph{Recall:} say the French for the cardboard box, then the French for the registered letter, then say \emph{la falaise} again
\end{itemize}

\begin{tcolorbox}[breakable,colback=teal!4,colframe=teal!35!black,title={Before you move on}]
What does la falaise mean? (\textquotedblleft{}The cliff\textquotedblright{}.) Say it once more.
\end{tcolorbox}
\section[le plateau]{le plateau — the plateau, the high plain; the tray}

\label{lesson:FR-C454-plateau}



\begin{quote}
Before the new one: say the French for the registered letter, then the French for the cliff.
\end{quote}

\subsection*{You'll want to know: le plateau}
\textbf{le plateau} — \textquotedblleft{}the plateau, the high plain; the tray\textquotedblright{}.

\begin{quote}
Le village est sur un plateau, au-dessus de la vallée.
\end{quote}

\begin{quote}
\emph{The village is on a plateau above the valley.}
\end{quote}

Plural \emph{les plateaux}. \emph{Un plateau} is also a tray, and a TV set.

\begin{cousinweb}
\emph{plateau} is built on \emph{plat} 'flat', which goes back to Vulgar Latin \emph{plattus}, from Greek \emph{platýs} 'broad' (EtymDB).
\end{cousinweb}

\begin{grammarlens}[title={What you've built}]
One more word for telling, explaining and describing at B1.
\end{grammarlens}

\subsection*{Your turn}
\begin{itemize}
\raggedright
  \item \emph{Say it:} \emph{le plateau}
  \item \emph{Say it:} \emph{le plateau}, once more
  \item \emph{Say it:} a sentence of your own with \emph{le plateau}
  \item \emph{Recall:} say the French for the registered letter, then the French for the cliff, then say \emph{le plateau} again
\end{itemize}

\begin{tcolorbox}[breakable,colback=teal!4,colframe=teal!35!black,title={Before you move on}]
What does le plateau mean? (\textquotedblleft{}The plateau, the high plain; the tray\textquotedblright{}.) Say it once more.
\end{tcolorbox}
\section[le sommet]{le sommet — the summit, the top; the summit meeting}

\label{lesson:FR-C454-sommet}



\begin{quote}
Before the new one: say the French for the cliff, then the French for the plateau.
\end{quote}

\subsection*{You'll want to know: le sommet}
\textbf{le sommet} — \textquotedblleft{}the summit, the top; the summit meeting\textquotedblright{}.

\begin{quote}
Nous sommes arrivés au sommet à midi.
\end{quote}

\begin{quote}
\emph{We reached the summit at noon.}
\end{quote}

\emph{Un sommet européen} — a European summit.

\begin{cousinweb}
\emph{sommet} comes from Old French \emph{somet} (EtymDB), a small \emph{som} 'top', from Latin \emph{summum} 'highest point', from \emph{summus} 'highest'.
\end{cousinweb}

\begin{grammarlens}[title={What you've built}]
One more word for telling, explaining and describing at B1.
\end{grammarlens}

\subsection*{Your turn}
\begin{itemize}
\raggedright
  \item \emph{Say it:} \emph{le sommet}
  \item \emph{Say it:} \emph{le sommet}, once more
  \item \emph{Say it:} a sentence of your own with \emph{le sommet}
  \item \emph{Recall:} say the French for the cliff, then the French for the plateau, then say \emph{le sommet} again
\end{itemize}

\begin{tcolorbox}[breakable,colback=teal!4,colframe=teal!35!black,title={Before you move on}]
What does le sommet mean? (\textquotedblleft{}The summit, the top; the summit meeting\textquotedblright{}.) Say it once more.
\end{tcolorbox}
\section[le col]{le col — the mountain pass; the collar}

\label{lesson:FR-C454-col}



\begin{quote}
Before the new one: say the French for the plateau, then the French for the summit.
\end{quote}

\subsection*{You'll want to know: le col}
\textbf{le col} — \textquotedblleft{}the mountain pass; the collar\textquotedblright{}.

\begin{quote}
La route du col est fermée en hiver.
\end{quote}

\begin{quote}
\emph{The road over the pass is closed in winter.}
\end{quote}

\emph{Le col de la chemise} — the shirt collar.

\begin{cousinweb}
\emph{col} is inherited through Old French \emph{col} 'neck' from Latin \emph{collum} 'neck' (EtymDB): a low point between two mountains, like a neck.
\end{cousinweb}

\begin{grammarlens}[title={What you've built}]
One more word for telling, explaining and describing at B1.
\end{grammarlens}

\subsection*{Your turn}
\begin{itemize}
\raggedright
  \item \emph{Say it:} \emph{le col}
  \item \emph{Say it:} \emph{le col}, once more
  \item \emph{Say it:} a sentence of your own with \emph{le col}
  \item \emph{Recall:} say the French for the plateau, then the French for the summit, then say \emph{le col} again
\end{itemize}

\begin{tcolorbox}[breakable,colback=teal!4,colframe=teal!35!black,title={Before you move on}]
What does le col mean? (\textquotedblleft{}The mountain pass; the collar\textquotedblright{}.) Say it once more.
\end{tcolorbox}
\section[le versant]{le versant — the mountainside, the slope}

\label{lesson:FR-C454-versant}



\begin{quote}
Before the new one: say the French for the summit, then the French for the mountain pass.
\end{quote}

\subsection*{You'll want to know: le versant}
\textbf{le versant} — \textquotedblleft{}the mountainside, the slope\textquotedblright{}.

\begin{quote}
Les vignes poussent sur le versant sud.
\end{quote}

\begin{quote}
\emph{The vines grow on the south-facing side.}
\end{quote}

\emph{Le versant nord} — the north face.

\begin{cousinweb}
\emph{versant} is the present participle of \emph{verser} 'to pour, to tip', used as a noun (EtymDB lists it); \emph{verser} goes back to Latin \emph{versāre}, from \emph{vertere} 'to turn' (Lewis \& Short).
\end{cousinweb}

\begin{grammarlens}[title={What you've built}]
One more word for telling, explaining and describing at B1.
\end{grammarlens}

\subsection*{Your turn}
\begin{itemize}
\raggedright
  \item \emph{Say it:} \emph{le versant}
  \item \emph{Say it:} \emph{le versant}, once more
  \item \emph{Say it:} a sentence of your own with \emph{le versant}
  \item \emph{Recall:} say the French for the summit, then the French for the mountain pass, then say \emph{le versant} again
\end{itemize}

\begin{tcolorbox}[breakable,colback=teal!4,colframe=teal!35!black,title={Before you move on}]
What does le versant mean? (\textquotedblleft{}The mountainside, the slope\textquotedblright{}.) Say it once more.
\end{tcolorbox}
